
\documentclass{article}
\usepackage{colortbl}
\usepackage{makecell}
\usepackage{multirow}
\usepackage{supertabular}

\begin{document}

\newcounter{utterance}

\centering \large Interaction Transcript for game `taboo', experiment `high\_en', episode 17 with reflexion{-}Qwen3.5{-}4B.
\vspace{24pt}

{ \footnotesize  \setcounter{utterance}{1}
\setlength{\tabcolsep}{0pt}
\begin{supertabular}{c@{$\;$}|p{.15\linewidth}@{}p{.15\linewidth}p{.15\linewidth}p{.15\linewidth}p{.15\linewidth}p{.15\linewidth}}
    \# & $\;$A & \multicolumn{4}{c}{Game Master} & $\;\:$B\\
    \hline

    \theutterance \stepcounter{utterance}  
    & & \multicolumn{4}{p{0.6\linewidth}}{
        \cellcolor[rgb]{0.9,0.9,0.9}{
            \makecell[{{p{\linewidth}}}]{
                \texttt{\tiny{[P1$\langle$GM]}}
                \texttt{You are playing a collaborative word guessing game in which you have to describe a target word for another player to guess.} \\
\\ 
\texttt{Rules:} \\
\texttt{(a) You have to reply in the form: CLUE: <some text>. Guesses from the other player will start with GUESS.} \\
\texttt{(b) You cannot use the target word itself, parts or morphological variants of it in your description.} \\
\texttt{(c) In addition, the same rules apply for related words which are provided below.} \\
\\ 
\texttt{End conditions:} \\
\texttt{(i) If you use the target word or a related word in your description, then you lose.} \\
\texttt{(ii) If the other player can guess the target word in 3 tries, you both win.} \\
\\ 
\texttt{Let us start.} \\
\\ 
\texttt{This is the target word that you need to describe and that the other player needs to guess:} \\
\\ 
\texttt{pricing} \\
\\ 
\texttt{Related words are:} \\
\\ 
\texttt{{-} evaluation} \\
\texttt{{-} cost} \\
\texttt{{-} value} \\
\\ 
\texttt{Important: You are under time pressure, give short descriptions that are to the point!} \\
            }
        }
    }
    & \\ \\

    \theutterance \stepcounter{utterance}  
    & \multicolumn{4}{p{0.6\linewidth}}{
        \cellcolor[rgb]{0.9,0.9,0.9}{
            \makecell[{{p{\linewidth}}}]{
                \texttt{\tiny{[P1$\rangle$GM]}}
                \texttt{CLUE: The amount of money a seller asks for a product or service.} \\
            }
        }
    }
    & & \\ \\

    \theutterance \stepcounter{utterance}  
    & & & \multicolumn{2}{p{0.3\linewidth}}{
        \cellcolor[rgb]{0.9,0.9,0.9}{
            \makecell[{{p{\linewidth}}}]{
                \texttt{\tiny{[GM$|$GM]}}
                \texttt{continue} \\
            }
        }
    }
    & & \\ \\

    \theutterance \stepcounter{utterance}  
    & & & \multicolumn{2}{p{0.3\linewidth}}{
        \cellcolor[rgb]{0.9,0.9,0.9}{
            \makecell[{{p{\linewidth}}}]{
                \texttt{\tiny{[GM$|$GM]}}
                \texttt{the amount of money a seller asks for a product or service} \\
            }
        }
    }
    & & \\ \\

    \theutterance \stepcounter{utterance}  
    & & \multicolumn{4}{p{0.6\linewidth}}{
        \cellcolor[rgb]{0.9,0.9,0.9}{
            \makecell[{{p{\linewidth}}}]{
                \texttt{\tiny{[GM$\rangle$P2]}}
                \texttt{You are playing a collaborative word guessing game in which you have to guess a target word that another player describes to you.} \\
\\ 
\texttt{You can make one guess at each trial. You win when you guess the target word. You lose when you cannot guess it in 3 tries.} \\
\\ 
\texttt{After each trial you will get a new hint from the other player which starts with CLUE.} \\
\\ 
\texttt{Make your guesses by just saying the word using the following form: GUESS: <a word>} \\
\\ 
\texttt{Let us start.} \\
\\ 
\\ 
\texttt{CLUE: The amount of money a seller asks for a product or service.} \\
            }
        }
    }
    & \\ \\

    \theutterance \stepcounter{utterance}  
    & & & \multicolumn{4}{p{0.6\linewidth}}{
        \cellcolor[rgb]{0.9,0.9,0.9}{
            \makecell[{{p{\linewidth}}}]{
                \texttt{\tiny{[GM$\langle$P2]}}
                \texttt{GUESS: price} \\
            }
        }
    }
     \\ \\

    \theutterance \stepcounter{utterance}  
    & & & \multicolumn{2}{p{0.3\linewidth}}{
        \cellcolor[rgb]{0.9,0.9,0.9}{
            \makecell[{{p{\linewidth}}}]{
                \texttt{\tiny{[GM$|$GM]}}
                \texttt{continue} \\
            }
        }
    }
    & & \\ \\

    \theutterance \stepcounter{utterance}  
    & & & \multicolumn{2}{p{0.3\linewidth}}{
        \cellcolor[rgb]{0.9,0.9,0.9}{
            \makecell[{{p{\linewidth}}}]{
                \texttt{\tiny{[GM$|$GM]}}
                \texttt{price} \\
            }
        }
    }
    & & \\ \\

    \theutterance \stepcounter{utterance}  
    & & \multicolumn{4}{p{0.6\linewidth}}{
        \cellcolor[rgb]{0.9,0.9,0.9}{
            \makecell[{{p{\linewidth}}}]{
                \texttt{\tiny{[P1$\langle$GM]}}
                \texttt{GUESS: price} \\
            }
        }
    }
    & \\ \\

    \theutterance \stepcounter{utterance}  
    & \multicolumn{4}{p{0.6\linewidth}}{
        \cellcolor[rgb]{0.9,0.9,0.9}{
            \makecell[{{p{\linewidth}}}]{
                \texttt{\tiny{[P1$\rangle$GM]}}
                \texttt{CLUE: That's close, but "price" is the noun for the specific amount, while the target word is the broader concept of setting that amount. Think of it as the act of determining the cost.} \\
            }
        }
    }
    & & \\ \\

    \theutterance \stepcounter{utterance}  
    & & & \multicolumn{2}{p{0.3\linewidth}}{
        \cellcolor[rgb]{0.9,0.9,0.9}{
            \makecell[{{p{\linewidth}}}]{
                \texttt{\tiny{[GM$|$GM]}}
                \texttt{continue} \\
            }
        }
    }
    & & \\ \\

    \theutterance \stepcounter{utterance}  
    & & & \multicolumn{2}{p{0.3\linewidth}}{
        \cellcolor[rgb]{0.9,0.9,0.9}{
            \makecell[{{p{\linewidth}}}]{
                \texttt{\tiny{[GM$|$GM]}}
                \texttt{Target word 'pricing' (stem=price) is similar to clue word 'price' (stem=price)} \\
            }
        }
    }
    & & \\ \\

    \theutterance \stepcounter{utterance}  
    & \multicolumn{4}{p{0.6\linewidth}}{
        \cellcolor[rgb]{0.95,0.95,0.95}{
            \makecell[{{p{\linewidth}}}]{
                \texttt{\tiny{[P2$\rightarrow$P1]}}
                \texttt{Post{-}episode Reflexion (not sent to GM):} \\
\texttt{New plan: I failed because my second clue was too abstract ("act of determining the cost") and didn't provide a concrete noun that the other player could easily guess within three tries. The target word "pricing" is a gerund/noun, so I should have used a clue that directly points to the action of setting a price or the concept of a price list, rather than a vague definition. Next time, I will use a specific, concrete noun phrase as my clue, such as "the act of setting a price" or "establishing a price," ensuring it is distinct from "price" itself but clearly leads to "pricing." I will also avoid over{-}explaining the difference between "price" and "pricing" in the clue itself, as that confuses the guesser.} \\
            }
        }
    }
    & & \\ \\

    \theutterance \stepcounter{utterance}  
    & & & \multicolumn{4}{p{0.6\linewidth}}{
        \cellcolor[rgb]{0.95,0.95,0.95}{
            \makecell[{{p{\linewidth}}}]{
                \texttt{\tiny{[P1$\rightarrow$P2]}}
                \texttt{Post{-}episode Reflexion (not sent to GM):} \\
\texttt{New plan: I guessed "price" immediately, which was likely too specific and narrow. The target word is a synonym for price, but "price" itself might not be the answer if the clue describes the concept generally. Next time, I should start with a broader, more common synonym like "cost" or "value" to increase the probability of hitting the target on the first or second try, rather than risking a loss with a single guess.} \\
            }
        }
    }
     \\ \\

\end{supertabular}
}

\end{document}
